\documentclass[a4paper,12pt]{article}
\usepackage[T2A]{fontenc}
\usepackage[utf8]{inputenc}
\usepackage[russian]{babel}
\usepackage[usenames,dvipsnames]{color}
\usepackage[version=4]{mhchem}
\usepackage[
  a4paper,
  top = 1.50cm,
  bottom = 1.50cm,
  left = 1.70cm,
  right = 1.70cm
]{geometry}
\usepackage{fancyhdr}
\usepackage{cmbright}
\usepackage{hyperref}
\hypersetup{
    colorlinks=true,
    linkcolor=RoyalBlue,
    urlcolor=RoyalBlue,
    citecolor=RoyalBlue,
    anchorcolor=RoyalBlue
}
\urlstyle{same}
\input{glyphtounicode}
\renewcommand{\baselinestretch}{1.05}
\renewcommand{\headrulewidth}{0pt}
\renewcommand{\footrulewidth}{0pt}
\setlength{\footskip}{30pt}
\pagestyle{fancy}
\fancyhf{}
\raggedbottom
\setlength{\tabcolsep}{0in}
\begin{document}
\begin{center}
    {\Huge{\fontsize{25}{30}\selectfont{Реза}} {\fontsize{25}{30}\selectfont{\textbf{Алиасгари Ренани}}}}
    \\
    {\small \raisebox{-0.2\height}\ {17/11/2025}} $|$
    {\small \raisebox{-0.2\height}\ {Реферат}} $|$
    {\small \raisebox{-0.2\height}\ {Российская Федерация}}
    {\small \raisebox{-0.2\height}\ {Аспирантура}}
    \\
    {\small \raisebox{-0.2\height}\ {Системный анализ, управление и обработка информации, статистика}}
    \vspace{-10pt}
\end{center}
\vspace{-10pt}
\noindent\rule{\textwidth}{0.5pt}

Моя исследовательская цель — продвинуть область электронной инженерии путём разработки радиационно-стойких видеопроцессоров на базе FPGA для космических применений, сосредоточившись на оптимизации алгоритмов обработки сигналов изображений для аппаратной реализации. Программа аспирантуры по специальности «Системный анализ, управление и обработка информации, статистика» в Российской Федерации и конкретно в Московском физико-техническом институте (МФТИ) предоставляет идеальную платформу для достижения этой цели, опираясь напрямую на мои предыдущие исследования и обучение в университете.
\vspace{10pt}

Я окончил бакалавриат по технической физике и магистратуру по прикладной математике и физике в МФТИ, где моя работа по изучению эффектов радиации в полупроводниках и разработке программируемых логических интегральных схем (FPGA) углубила мой интерес к созданию радиационно-стойких аппаратных систем. Я хочу разрабатывать видеопроцессоры на базе FPGA и Application-Specific Integrated Circuit (ASIC), интегрируя техники системного анализа и управления для повышения пропускной способности, снижения задержек и повышения надёжности в суровых условиях.
\vspace{10pt}

Более двух лет назад я присоединился к лаборатории доктора Ковешникова в Институте проблем технологии микроэлектроники и особочистых материалов Российской академии наук для изучения эффектов радиации на микроэлектронных структурах, сосредоточившись на электрической характеризации полупроводников и устройств Metal-Oxide-Semiconductor (MOS) на базе \ce{SiO2}. Мы изучали влияние электронной и гамма-иррадиации, а также обработки водородной плазмой на MOS-устройства, идентифицировали ловушки по местоположению (оксид, полупроводник, интерфейс) и типу носителей (электроны, дырки). Эта работа привела к публикации\footnotemark[1] с первым авторством и дала понимание стабильности устройств под напряжением смещения, иммунитета к радиации и плотности ловушек в оксиде на интерфейсе диэлектрик-полупроводник. Я также разработал оптимизированные приложения MATLAB для автоматизации экспериментальных методов характеризации, что сократило время выполнения и создало конвейеры измерений. Помимо углубления моего понимания физики полупроводников, это исследование также дало мне неоценимый опыт работы с экспериментальным оборудованием и автоматизации экспериментальных техник с использованием MATLAB и Python.
\vspace{10pt}

После окончания бакалавриата я присоединился к лаборатории разработки систем на кристалле (дизайн-центр по проектированию микропроцессорной техники для систем с искусственным интеллектом) и лаборатории плазменных систем под руководством доктора Васильева и доктора Чеснокова для разработки систем на базе FPGA и изучения эффектов радиации и плазмы на этих структурах. Я реализовал алгоритмы обработки видео путём ручного кодирования на Verilog и генерации HDL из моделей Simulink. Я верифицировал функциональность алгоритмов с использованием тестовых стендов и скриптов автоматизации на Python. В лаборатории плазмы мы проводили эксперименты, где FPGA с прикреплённой камерой и синтезированным дизайном обработки изображений помещалась в систему электронно-лучевой плазмы с низким давлением кислорода и подвергалась электронной иррадиации и рентгеновским лучам, генерируемым от вольфрамовой мишени. Сигнал выхода FPGA мониторился во время термоциклирования, и проводились предварительные тесты функциональности.
\vspace{10pt}

\newpage

В лаборатории разработки систем на кристалле, которая работает под компанией YADRO, я портировал математические алгоритмы в эффективные реализации на Verilog, построил библиотеку фиксированной точки и создал аппроксимации функций на базе LUT с использованием метода Горнера для поддержки вычислений с фиксированной точкой. Я реализовал ряд алгоритмов обработки изображений, включая преобразование rgb2hsv, сегментацию цвета, детекцию краёв Собеля, глобальное тональное отображение, демозаику, преобразование вида с высоты птичьего полёта, коррекцию рыбьего глаза и ресайзер изображений, через генерацию HDL-кода Simulink и рукописный Verilog. Я оптимизировал их по задержке и пропускной способности, решая проблемы синхронизации и конвейеризации.

\vspace{10pt}

Для верификации я разработал всесторонние тестовые стенды на Verilog и использовал Python для автоматизации симуляции и анализа данных, чтобы подтвердить функциональность и производительность DSP. Затем я синтезировал, отобразил и маршрутизировал HDL-код с использованием Vivado на платах Xilinx Zybo Z7, верифицируя алгоритмы ISP сначала путём потоковой передачи тестовых изображений через HDMI с хост-компьютера, а позже с живой камерой, подключённой напрямую к FPGA. Этот практический опыт в реализации и тестировании аппаратного обеспечения укрепил моё сотрудничество с YADRO.

\vspace{10pt}

Мой опыт работы с облучёнными полупроводниками, плазменными системами и дизайном FPGA вдохновил меня на продолжение аспирантуры в МФТИ по специальности «Системный анализ, управление и обработка информации, статистика». Я планирую исследовать эффекты радиации на устройствах FPGA, идентифицируя воздействия на компоненты, такие как RAM и MOSFET, и проектировать более стойкие архитектуры, такие как замена SRAM на ReRAM для повышения толерантности в космических средах. Это согласуется с моей текущей работой в дизайн-центре по проектированию микропроцессорной техники для систем с искусственным интеллектом и лаборатории разработки систем на кристалле в МФТИ, где я могу продолжить оптимизацию алгоритмов ISP для видеопроцессоров. Моя работа по генерации HDL-кода для алгоритмов обработки видео грядущая и будет представлена\footnotemark[2] на конференции FPGA-Systems 2025 в Москве. Я с нетерпением жду продолжения моих исследований в России для развёртывания этих систем для обработки сигналов в реальном времени в аэрокосмических приложениях.

\vspace{10pt}

Как иностранный студент, уже интегрированный в исследовательское сообщество МФТИ, включая сотрудничество с Институтом проблем технологии микроэлектроники и особочистых материалов Российской академии наук, я хорошо подготовлен к вкладу и получению пользы от его совместной среды. Меня мотивируют проекты, сочетающие экспериментальное тестирование с оптимизацией аппаратного обеспечения, и я с нетерпением жду применения моих навыков в Verilog и анализе радиации для продвижения видеопроцессоров на базе FPGA, потенциально поддерживая космические инициативы России. Для меня будет привилегией продолжить образование в МФТИ в России, центре научных инноваций с сильным акцентом на практические инженерные решения.

\footnotetext[1]{Р. Алиасгари Ренани, О.А. Солтанович, М.А. Князев, С.В. Ковешников, Исследование устройств MOS на основе \ce{SiO2}, облученных низкоэнергетичными электронами, методами емкость-напряжение и термостимулированный ток. \href{https://doi.org/10.1134/S1063739723600516}{Журнальная статья}, Russian Microelectronics, 2023}
\footnotetext[2]{В. Вологин, Р. Алиасгари Ренани, В. Васильевский, В. Чесноков, Сравнительный анализ ручного Verilog и HDL-кода, сгенерированного в Simulink, для алгоритмов обработки изображений. Предстоящая презентация, \href{https://meetups.yadro.com/fpga-msk-1125/}{Конференция FPGA-Systems 2025}}

\end{document}